\documentclass[10pt,a4paper]{amsart}
\usepackage{amsmath,amssymb,amsthm,graphicx}
\usepackage{verbatim}
\newenvironment{reference}{\comment}{\endcomment}
%\newenvironment{reference}{}{}
\newenvironment{slogan}{\comment}{\endcomment}
\newenvironment{history}{\comment}{\endcomment}

% For commutative diagrams we use Xy-pic
\usepackage[all]{xy}

% We use 2cell for 2-commutative diagrams.
\xyoption{2cell}
% Theorem environments.
%
\theoremstyle{plain}
\newtheorem{theorem}[subsection]{Theorem}
\newtheorem{proposition}[subsection]{Proposition}
\newtheorem{lemma}[subsection]{Lemma}

\theoremstyle{definition}
\newtheorem{definition}[subsection]{Definition}
\newtheorem{example}[subsection]{Example}
\newtheorem{exercise}[subsection]{Exercise}
\newtheorem{situation}[subsection]{Situation}

\theoremstyle{remark}
\newtheorem{remark}[subsection]{Remark}
\newtheorem{remarks}[subsection]{Remarks}

\numberwithin{equation}{subsection}

% Macros
%
\def\lim{\mathop{\mathrm{lim}}\nolimits}
\def\colim{\mathop{\mathrm{colim}}\nolimits}
\def\Spec{\mathop{\mathrm{Spec}}}
\def\Hom{\mathop{\mathrm{Hom}}\nolimits}
\def\Ext{\mathop{\mathrm{Ext}}\nolimits}
\def\SheafHom{\mathop{\mathcal{H}\!\mathit{om}}\nolimits}
\def\SheafExt{\mathop{\mathcal{E}\!\mathit{xt}}\nolimits}
\def\Sch{\mathit{Sch}}
\def\Mor{\mathop{\mathrm{Mor}}\nolimits}
\def\Ob{\mathop{\mathrm{Ob}}\nolimits}
\def\Sh{\mathop{\mathit{Sh}}\nolimits}
\def\NL{\mathop{N\!L}\nolimits}
\def\CH{\mathop{\mathrm{CH}}\nolimits}
\def\proetale{{pro\text{-}\acute{e}tale}}
\def\etale{{\acute{e}tale}}
\def\QCoh{\mathit{QCoh}}
\def\Ker{\mathop{\mathrm{Ker}}}
\def\Im{\mathop{\mathrm{Im}}}
\def\Coker{\mathop{\mathrm{Coker}}}
\def\Coim{\mathop{\mathrm{Coim}}}

% Boxtimes
%
\DeclareMathSymbol{\boxtimes}{\mathbin}{AMSa}{"02}

%
% Macros for moduli stacks/spaces
%
\def\QCohstack{\mathcal{QC}\!\mathit{oh}}
\def\Cohstack{\mathcal{C}\!\mathit{oh}}
\def\Spacesstack{\mathcal{S}\!\mathit{paces}}
\def\Quotfunctor{\mathrm{Quot}}
\def\Hilbfunctor{\mathrm{Hilb}}
\def\Curvesstack{\mathcal{C}\!\mathit{urves}}
\def\Polarizedstack{\mathcal{P}\!\mathit{olarized}}
\def\Complexesstack{\mathcal{C}\!\mathit{omplexes}}
% \Pic is the operator that assigns to X its picard group, usage \Pic(X)
% \Picardstack_{X/B} denotes the Picard stack of X over B
% \Picardfunctor_{X/B} denotes the Picard functor of X over B
\def\Pic{\mathop{\mathrm{Pic}}\nolimits}
\def\Picardstack{\mathcal{P}\!\mathit{ic}}
\def\Picardfunctor{\mathrm{Pic}}
\def\Deformationcategory{\mathcal{D}\!\mathit{ef}}
\usepackage{cleveref}
\crefname{originaldefinition}{original definition}{original definitions}
\Crefname{originaldefinition}{Original definition}{Original definitions}
\makeatletter
\newlabel{sites-modules-definition-2-morphism-ringed-topoi}{{04IC}{source}{}{sourcepointer-1}{}}
\newlabel{sites-modules-definition-2-morphism-ringed-topoi@cref}{{[originaldefinition][0][]04IC}{[1][0][]source}{}{}{}}
\makeatother
\crefname{originallemma}{original lemma}{original lemmas}
\Crefname{originallemma}{Original lemma}{Original lemmas}
\makeatletter
\newlabel{morphisms-lemma-quasi-affine-finite-type-over-S}{{04II}{source}{}{sourcepointer-2}{}}
\newlabel{morphisms-lemma-quasi-affine-finite-type-over-S@cref}{{[originallemma][0][]04II}{[1][0][]source}{}{}{}}
\makeatother
\crefname{originallemma}{original lemma}{original lemmas}
\Crefname{originallemma}{Original lemma}{Original lemmas}
\makeatletter
\newlabel{schemes-lemma-base-change-immersion}{{01JY}{source}{}{sourcepointer-3}{}}
\newlabel{schemes-lemma-base-change-immersion@cref}{{[originallemma][0][]01JY}{[1][0][]source}{}{}{}}
\makeatother
\crefname{originallemma}{original lemma}{original lemmas}
\Crefname{originallemma}{Original lemma}{Original lemmas}
\makeatletter
\newlabel{lemma-stalk-pullback}{{03Q1}{source}{}{sourcepointer-4}{}}
\newlabel{lemma-stalk-pullback@cref}{{[originallemma][0][]03Q1}{[1][0][]source}{}{}{}}
\makeatother
\crefname{originallemma}{original lemma}{original lemmas}
\Crefname{originallemma}{Original lemma}{Original lemmas}
\makeatletter
\newlabel{topologies-lemma-alternative}{{04HR}{source}{}{sourcepointer-5}{}}
\newlabel{topologies-lemma-alternative@cref}{{[originallemma][0][]04HR}{[1][0][]source}{}{}{}}
\makeatother
\crefname{originallemma}{original lemma}{original lemmas}
\Crefname{originallemma}{Original lemma}{Original lemmas}
\makeatletter
\newlabel{sites-lemma-sheaf-coequalizer-representable}{{00WS}{source}{}{sourcepointer-6}{}}
\newlabel{sites-lemma-sheaf-coequalizer-representable@cref}{{[originallemma][0][]00WS}{[1][0][]source}{}{}{}}
\makeatother
\title{The only 2-morphism of the locally ringed topos attached to a morphism of schemes is the identity}
\author{The Stacks Project Authors}
\date{stacks-project:etale-cohomology.tex:local-snapshot-20261011; GNU FDL 1.2 or later}
\begin{document}
\maketitle
\noindent\textit{Editorial note.} All delivered mathematical text is the author's own, copied from the pinned Stacks Project chapter file etale-cohomology.tex; no statement or proof step was rewritten, reordered or removed.\par\medskip
\noindent\textit{Editorial note.} The proof is the author's complete argument: the reduction to representable sheaves, the injection test and the propagation to all sheaves are written out; the parenthetical that one commutativity verification is omitted concerns a diagram chase and is preserved verbatim in the counted span.\par\medskip
\noindent\textit{Editorial note.} Every labelled reference occurring inside the counted spans is carried in assembly.external\_references as a non-load-bearing navigation pointer (same-chapter ones at their definition site, cross-chapter ones at the invocation inside the counted span), so no load-bearing mathematics is replaced by a stub. No bibliography citation occurs inside the counted spans.\par\medskip
\noindent\textit{Editorial note.} The checked chapter list for same-name collisions is the set of stacks chapters staged in this repository: chow, cohomology, divisors, etale-cohomology, flat, more-algebra, more-morphisms, morphisms, spaces-more-morphisms, varieties together with this card's chapter etale-cohomology.\par\medskip
\noindent\textit{Editorial note.} The preamble is the chapter's own native preamble excerpt already used for this source: the Xy-pic setup (native preamble lines 14-18) and the native theorem declarations with the mathematical macros (native preamble lines 154-217).\par\medskip
\noindent\textit{Allowed standard prerequisites (editorial).}\begin{itemize}
\item 2-morphisms of ringed topoi. Modules on Sites, Definition sites-modules-definition-2-morphism-ringed-topoi. Supplies the data that the transformation must respect, namely the compatibility with the structure sheaf maps.
\item Embedding an affine etale scheme into affine space. Morphisms, Lemma morphisms-lemma-quasi-affine-finite-type-over-S, applied after choosing an affine etale neighbourhood. Produces injections of sheaves with which the transformation is tested.
\item Pullback of representable sheaves and generation by representables. Same chapter, Lemma lemma-stalk-pullback, together with Topologies, Lemma topologies-lemma-alternative and Sites, Lemma sites-lemma-sheaf-coequalizer-representable. Identifies the pullback of a representable sheaf and lets the identity on generators propagate to all sheaves.
\end{itemize}\medskip
\ifdefined\hypertarget\hypertarget{sourcepointer-1}{}\fi
\noindent\textit{Source reference.} Original definition 04IC: sources/stacks-etale-cohomology/pinned-original-source.tex, source lines 5231--5231. Navigation pointer invoked at proof line 5231; the counted span contains the pointer only and the cited result is not reproduced inside it, so the pointer is registered as a non-load-bearing navigation stub in the established form for Stacks material. No bibliography citation is involved.\par\medskip
\ifdefined\hypertarget\hypertarget{sourcepointer-2}{}\fi
\noindent\textit{Source reference.} Original lemma 04II: sources/stacks-etale-cohomology/pinned-original-source.tex, source lines 5247--5247. Navigation pointer invoked at proof line 5247; the counted span contains the pointer only and the cited result is not reproduced inside it, so the pointer is registered as a non-load-bearing navigation stub in the established form for Stacks material. No bibliography citation is involved.\par\medskip
\ifdefined\hypertarget\hypertarget{sourcepointer-3}{}\fi
\noindent\textit{Source reference.} Original lemma 01JY: sources/stacks-etale-cohomology/pinned-original-source.tex, source lines 5252--5252. Navigation pointer invoked at proof line 5252; the counted span contains the pointer only and the cited result is not reproduced inside it, so the pointer is registered as a non-load-bearing navigation stub in the established form for Stacks material. No bibliography citation is involved.\par\medskip
\ifdefined\hypertarget\hypertarget{sourcepointer-4}{}\fi
\noindent\textit{Source reference.} Original lemma 03Q1: sources/stacks-etale-cohomology/pinned-original-source.tex, source lines 4649--4649. Navigation pointer invoked at proof line 5204; the counted span contains the pointer only and the cited result is not reproduced inside it, so the pointer is registered as a non-load-bearing navigation stub in the established form for Stacks material. No bibliography citation is involved.\par\medskip
\ifdefined\hypertarget\hypertarget{sourcepointer-5}{}\fi
\noindent\textit{Source reference.} Original lemma 04HR: sources/stacks-etale-cohomology/pinned-original-source.tex, source lines 5290--5290. Navigation pointer invoked at proof line 5290; the counted span contains the pointer only and the cited result is not reproduced inside it, so the pointer is registered as a non-load-bearing navigation stub in the established form for Stacks material. No bibliography citation is involved.\par\medskip
\ifdefined\hypertarget\hypertarget{sourcepointer-6}{}\fi
\noindent\textit{Source reference.} Original lemma 00WS: sources/stacks-etale-cohomology/pinned-original-source.tex, source lines 5292--5292. Navigation pointer invoked at proof line 5292; the counted span contains the pointer only and the cited result is not reproduced inside it, so the pointer is registered as a non-load-bearing navigation stub in the established form for Stacks material. No bibliography citation is involved.\par\medskip
\begin{lemma}
\label{lemma-2-morphism}
Let $X$, $Y$ be schemes. Let $f : X \to Y$ be a morphism of schemes.
Let $t$ be a $2$-morphism from $(f_{small}, f_{small}^\sharp)$ to itself, see
Modules on Sites,
Definition \ref{sites-modules-definition-2-morphism-ringed-topoi}.
Then $t = \text{id}$.
\end{lemma}

\begin{proof}
This means that $t : f^{-1}_{small} \to f^{-1}_{small}$
is a transformation of functors such that the diagram
$$
\xymatrix{
f_{small}^{-1}\mathcal{O}_Y
\ar[rd]_{f_{small}^\sharp}  & &
f_{small}^{-1}\mathcal{O}_Y \ar[ll]^t \ar[ld]^{f_{small}^\sharp} \\
& \mathcal{O}_X
}
$$
is commutative. Suppose $V \to Y$ is \'etale with $V$ affine. By
Morphisms, Lemma \ref{morphisms-lemma-quasi-affine-finite-type-over-S}
we may choose an immersion $i : V \to \mathbf{A}^n_Y$ over $Y$.
In terms of sheaves this means that $i$ induces an injection
$h_i : h_V \to \prod_{j = 1, \ldots, n} \mathcal{O}_Y$ of sheaves.
The base change $i'$ of $i$ to $X$ is an immersion
(Schemes, Lemma \ref{schemes-lemma-base-change-immersion}).
Hence $i' : X \times_Y V \to \mathbf{A}^n_X$ is an immersion, which
in turn means that
$h_{i'} : h_{X \times_Y V} \to \prod_{j = 1, \ldots, n} \mathcal{O}_X$
is an injection of sheaves.
Via the identification $f_{small}^{-1}h_V = h_{X \times_Y V}$ of
Lemma \ref{lemma-stalk-pullback}
the map $h_{i'}$ is equal to
$$
\xymatrix{
f_{small}^{-1}h_V \ar[r]^-{f^{-1}h_i} &
\prod_{j = 1, \ldots, n} f_{small}^{-1}\mathcal{O}_Y
\ar[r]^{\prod f^\sharp} &
\prod_{j = 1, \ldots, n} \mathcal{O}_X
}
$$
(verification omitted). This means that the map
$t : f_{small}^{-1}h_V \to f_{small}^{-1}h_V$
fits into the commutative diagram
$$
\xymatrix{
f_{small}^{-1}h_V \ar[r]^-{f^{-1}h_i} \ar[d]^t &
\prod_{j = 1, \ldots, n} f_{small}^{-1}\mathcal{O}_Y
\ar[r]^-{\prod f^\sharp} \ar[d]^{\prod t} &
\prod_{j = 1, \ldots, n} \mathcal{O}_X \ar[d]^{\text{id}}\\
f_{small}^{-1}h_V \ar[r]^-{f^{-1}h_i} &
\prod_{j = 1, \ldots, n} f_{small}^{-1}\mathcal{O}_Y
\ar[r]^-{\prod f^\sharp} &
\prod_{j = 1, \ldots, n} \mathcal{O}_X
}
$$
The commutativity of the right square holds by our assumption on $t$
explained above.
Since the composition of the horizontal arrows is injective
by the discussion above we conclude that the left vertical arrow
is the identity map as well. Any sheaf of sets on
$Y_\etale$ admits a surjection from a (huge) coproduct of sheaves
of the form $h_V$ with $V$ affine (combine
Topologies, Lemma \ref{topologies-lemma-alternative}
with
Sites, Lemma \ref{sites-lemma-sheaf-coequalizer-representable}).
Thus we conclude that $t : f_{small}^{-1} \to f_{small}^{-1}$
is the identity transformation as desired.
\end{proof}

\end{document}
